\section*{第 \ref{ch:g12:climate-change} 章——气候变化：观测、成因与预测}

\begin{solution}{exo:g12:climate-change:1}
它所引起的、进入气候系统的净能量流的变化，
发生在大气层顶部，且早于地表温度的响应；
以 \unit{W/m^2} 计量。正值：行星获得能量而变暖；
负值：行星失去能量而变冷。
\end{solution}

\begin{solution}{exo:g12:climate-change:2}
五月最高，九月最低。春季和夏季，
北半球植被（大部分陆地植物生长于此）通过光合作用吸收的二氧化碳多于
植物和土壤释放的量；秋季和冬季，
分解与呼吸作用占主导，浓度又重新上升。
\end{solution}

\begin{solution}{exo:g12:climate-change:3}
\emph{很可能}：概率为 90 至 100\,\%；\emph{可能}：66 至
100\,\%。$[0.95 \text{ 至 } 1.20]$ 是一个很可能的范围：有 90\,\% 的把握认为
真实值位于 0.95 与 1.20 之间。
\end{solution}

\begin{solution}{exo:g12:climate-change:4}
91\,\%。水受热膨胀：变暖的海洋占据的体积更大，因此
即使不增加水量，其表面也会上升（陆地冰融化所增加的水量
会叠加在此之上）。
\end{solution}

\begin{solution}{exo:g12:climate-change:5}
水汽：变暖的空气能容纳更多水汽，水汽是一种温室气体，
会加剧变暖。雪与冰：融化后露出颜色更深的地表，
吸收更多阳光。负反馈：行星变暖后向太空辐射更多红外线，
使其重新回到平衡。
\end{solution}

\begin{solution}{exo:g12:climate-change:6}
$427.35 - 315.98 = \qty{111.37}{ppm}$；$111.37/315.98 \approx 0.352$，
增长约 35\,\%。66 年间：$111.37/66 \approx
\qty{1.7}{ppm}$ 每年平均。
\end{solution}

\begin{solution}{exo:g12:climate-change:7}
$F = 3.93 \times \ln(427.35/278)/\ln 2 = 3.93 \times 0.4300/0.6931 \approx
\qty{2.44}{W/m^2}$。
\end{solution}

\begin{solution}{exo:g12:climate-change:8}
$3.7 \times 30 = \qty{111}{mm}$，约为 \qty{11}{cm}。$3.7/1.3 \approx 2.8$：
速度几乎快了三倍。
\end{solution}

\begin{solution}{exo:g12:climate-change:9}
$\Delta T = 2.72/1.16 \approx \qty{2.3}{\degreeCelsius}$。气候尚未达到平衡：
海洋吸收了约 90\,\% 的额外能量，缓慢变暖，
因此地表尚未跟上强迫的变化（系统
仍在获得约 \qty{0.8}{W/m^2}）。
\end{solution}

\begin{solution}{exo:g12:climate-change:10}
太阳和火山至多只能解释 \qty{1.06}{\degreeCelsius} 变暖中的 \qty{0.1}{\degreeCelsius}，
它们不可能是其成因。另一种自然成因是
气候的内部变率（即其自然波动，例如
海洋与大气之间的热量交换）；评估给出
其贡献为 $-0.2$ 至 $+\qty{0.2}{\degreeCelsius}$，
同样太小。
\end{solution}

\begin{solution}{exo:g12:climate-change:11}
很可能低于 \qty{2}{\degreeCelsius}：只有 SSP1-1.9（很可能范围
为 1.0 至 1.8）。很可能高于：SSP2-4.5（2.1 至 3.5）、SSP3-7.0（2.8 至 4.6）
和 SSP5-8.5（3.3 至 5.7）。对于 SSP1-2.6（1.3 至 2.4），两者均不很可能。
\end{solution}

\begin{solution}{exo:g12:climate-change:12}
气溶胶强迫为 $-0.22 - 0.84 = \qty{-1.06}{W/m^2}$；若无此项，
总强迫变为 $2.72 + 1.06 = \qty{3.78}{W/m^2}$。平衡变暖
增加 $1.06/1.16 \approx \qty{0.9}{\degreeCelsius}$。
\end{solution}

\begin{solution}{exo:g12:climate-change:13}
由于 $\ln(2C) - \ln C = \ln 2$ 对任意 $C$ 成立，
\[
  F(556) - F(278) = \frac{3.93\,\ln 2}{\ln 2} = \qty{3.93}{W/m^2}, \qquad
  F(1112) - F(556) = \frac{3.93\,\ln 2}{\ln 2} = \qty{3.93}{W/m^2}:
\]
每次浓度翻倍所产生的强迫相同。在 \qty{1112}{ppm} 时的总强迫为 $2 \times 3.93 =
\qty{7.86}{W/m^2}$。
\end{solution}

\begin{solution}{exo:g12:climate-change:14}
$\dfrac{\dd F}{\dd C} = \dfrac{F_{2\times}}{C \ln 2}$。在 \qty{280}{ppm} 时：
$3.93/(280 \times 0.6931) \approx \qty{0.0202}{W/m^2}$ 每 ppm；在
\qty{420}{ppm} 时：$\approx \qty{0.0135}{W/m^2}$ 每 ppm。随着浓度增加，每增加一 ppm
所产生的强迫减小（在 420 时比在 280 时小 1.5 倍），
但浓度仍在持续上升，因此强迫依然在增长。
\end{solution}

\begin{solution}{exo:g12:climate-change:15}
无反馈时：$3.93/3.22 \approx \qty{1.2}{\degreeCelsius}$；
有反馈时：$3.93/1.16 \approx \qty{3.4}{\degreeCelsius}$。放大
因子：$3.22/1.16 \approx 2.8$。
\end{solution}

\begin{solution}{pb:g12:climate-change:1}
\textbf{1.} $3.37 + 3.51 = \qty{6.88}{ppm}$，约为 \qty{7}{ppm}。

\textbf{2.} 从春季到夏末，北半球的植物
通过光合作用从空气中吸收二氧化碳的速度快于
呼吸与分解的释放，浓度一直下降至九月。
秋季和冬季，二者的平衡逆转，浓度重新攀升，
回到五月的最高值。

\textbf{3.} 南半球陆地远少于北半球，因而
植被也少得多；其表面大部分为海洋。

\textbf{4.} 该采样点位置较高，位于近地面污染空气层之上；
远离大陆和城市，其采集的空气混合充分，
代表整个大气，而非局地源。

\textbf{5.} $(325.68 - 316.91)/10 \approx \qty{0.88}{ppm}$ 每年。

\textbf{6.} $(427.35 - 401.01)/10 \approx \qty{2.63}{ppm}$ 每年。

\textbf{7.} $2.634/0.877 \approx 3.0$：增长速率已增至三倍。

\textbf{8.} $(427.35 - 278)/278 \approx 0.537$：约为 54\,\%。

\textbf{9.} $427.35 - 300 \approx \qty{127}{ppm}$。

\textbf{10.} $F = 3.93 \times \ln(427.35/278)/\ln 2 \approx
\qty{2.44}{W/m^2}$。

\textbf{11.} $F(316.91) = 3.93 \times \ln(316.91/278)/\ln 2 \approx
\qty{0.74}{W/m^2}$；增量为 $2.44 - 0.74 \approx \qty{1.70}{W/m^2}$。

\textbf{12.} 再增加 \unit{W/m^2}，浓度乘以
$2^{1/3.93} \approx 1.193$：$427.35 \times 1.193 \approx \qty{510}{ppm}$。

\textbf{13.} $(509.8 - 427.35)/2.634 \approx 31$ 年。

\textbf{14.} $2 \times 278 = \qty{556}{ppm}$；\omterm{def:g12:climate-change:forcing}{辐射强迫} $F_{2\times} =
\qty{3.93}{W/m^2}$。

\textbf{15.} 由敏感度的定义，\qty{3}{\degreeCelsius}
（很可能处于 2.5 至 \qty{4}{\degreeCelsius} 之间）。

\textbf{16.} $\lambda = 3.93/3 \approx \qty{1.31}{W.m^{-2}.\degreeCelsius^{-1}}$。
反馈之和为 \qty{1.16}{W.m^{-2}.\degreeCelsius^{-1}}，由此得
$3.93/1.16 \approx \qty{3.4}{\degreeCelsius}$，略大一些；评估的
最佳估计综合了这一物理过程、观测到的变暖以及过去的
气候，两个数值均处于可能范围之内。

\textbf{17.} 海洋吸收了额外能量中的大部分，升温缓慢：
表层要经过数十年至数百年才能达到平衡增温。

\textbf{18.} 它假设未来排放的一种情景（过去十年的增长永远持续），
这就是一种\omterm{def:g12:climate-change:scenario}{排放情景}。\omterm{def:g12:climate-change:scenario}{预估}则借助物理模式，计算气候与碳循环
对某一\omterm{def:g12:climate-change:scenario}{排放情景}的响应，
并给出范围与可能性。

\textbf{19.} $(556 - 427.35)/2.634 \approx 48.8$ 年，即 2025 年之后约
\textbf{2074}。
\end{solution}
